
Soit $(x_i)_{i\in I}\in K_\bullet(S)^{(I)}$~; on a :
\[
\varphi_U\circ j^*((x_i))=\sum_i \xi_i\cdot f_U^*\circ j^*(x_i);
\]
\[
v^*\circ\varphi((x_i))=v^*\left(\sum_i \xi_i\cdot f^*(x_i)\right)=\sum_i \xi_i\cdot v^*\circ f^*(x_i)=\sum_i \xi_i\cdot f_U^*\circ j^*(x_i)
\]
l'égalité du centre résultant de IV 2.12.2.

Ce résultat permet de prouver 1.2 par récurrence noethérienne sur l'ensemble des sous-schémas fermés de $S$. Si on suppose que pour tout sous-schéma fermé $S'$ de $S$, distinct de $S$, $\varphi_{S'}$ est surjective, alors le raisonnement précédent montre qu'on peut se ramener à prouver 1.2 pour un schéma de base irréductible et affine, quitte à remplacer $S$ par un ouvert affine irréductible. De plus, le diagramme
\[
\begin{tikzcd}
K_\bullet(S_{\red})^{(I)} \arrow[r] \arrow[d,"\varphi_{S_{\red}}"'] & K_\bullet(S)^{(I)} \arrow[d,"\varphi"]\\
K_\bullet(X_{S_{\red}}) \arrow[r] & K_\bullet(X),
\end{tikzcd}
\]
où $K_\bullet(X_{S_{\red}})\to K_\bullet(X)$ est surjectif d'après 1.1, montre qu'on peut également supposer $S$ réduit.

On se place donc sur un schéma de base intègre et affine. On pose $S=\Spec(A)$, et on note $K$ le corps des fractions de $A$; si $s$ est le point générique de $S$, alors $\underline{O}_{S,s}=K$.

Soit $y\in K_\bullet(X)$. Par hypothèse, $K_\bullet(X_s)$ est engendré par les images des $\xi_i$; si $y_s$ est l'image de $y$ dans $K_\bullet(X_s)$, on peut trouver des entiers $n_i$ tels que :
\[
y_s=\sum_i\xi_i\cdot n_i.
\]
